\section{Result 1: with 8 bits per animal, no model ran out of room}\label{sec:round1}

This is the first round, without width 4: the four languages (16, 64, 256 and 1,024 animals),
three widths (8, 16 and 32) and three repeats of each pairing, 36 runs in all. In the trial
runs width 4 had still learned no facts after 360 sets, even on 16 animals, so I ran it
separately (section \ref{sec:four}). That is a loss for the test: by the rough guide width 4 was
the one size that would have come near full on 1,024 animals.

\begin{table}[H]
\centering\small
\begin{tabular}{@{}rrrrr@{}}
\toprule
 & \multicolumn{4}{c}{Animals in the vocabulary} \\
Width & 16 & 64 & 256 & 1{,}024 \\
\midrule
8 & 100 & 100 & 100 & 100 \\
16 & 100 & 100 & 100 & 100 \\
32 & 100 & 100 & 100 & 100 \\
\bottomrule
\end{tabular}

\caption{Facts known (0 to 100) at the end of each run, with 8 bits per animal. Each entry
covers three repeats; a single number means all three gave that score.}
\label{tab:round1}
\end{table}

\textbf{All 36 runs ended as solved, with a score of 100} (table \ref{tab:round1}). A model of
width 8 learned all 4,096 facts about 1,024 animals as completely as it learned the 64 facts
about 16 animals.

Remember the catch from section \ref{sec:model}: this was not one store of fixed size taking
in 64 times as much. Every added animal came with a row of its own, and only the body was the
same size in both.

\textbf{Verdict of this round: undecided}, as section \ref{sec:rule} said it would be if no
model did worse. The rule is not met for width 8 or for width 16, the two widths it could be
met for. That is not evidence against your idea, because the idea predicts a loss only once
the words no longer fit, and by the rough guide the fullest of these models (width 8 on 1,024
animals) was a little over one-third full. What the round does show is narrower: with room to
spare, 64 times as many words did not leave any of these models worse off at the end of its
training.

\textbf{What a larger vocabulary did cost was training time.} Table \ref{tab:effort1} shows
how many sets each run had read when it was declared solved.

\begin{table}[H]
\centering\small
\begin{tabular}{@{}rrrrr@{}}
\toprule
 & \multicolumn{4}{c}{Animals in the vocabulary} \\
Width & 16 & 64 & 256 & 1{,}024 \\
\midrule
8 & 64, 64, 91 & 64, 91, 181 & 256, 256, 512 & 256, 512, 1{,}639 \\
16 & 32, 32, 32 & 45, 45, 64 & 64, 64, 91 & 128, 181, 181 \\
32 & 16, 23, 23 & 23, 32, 32 & 64, 64, 64 & 91, 128, 128 \\
\bottomrule
\end{tabular}

\caption{Sets of training text read before each run was solved, for the three repeats,
smallest first. Runs are graded only after certain numbers of sets and must pass twice in a
row, so these figures move in jumps, and two equal figures may hide a real difference. The
figure 1,639 is the one run that was graded because its time had run out.}
\label{tab:effort1}
\end{table}

Two things stand out.

\textbf{More words took more sets, but far fewer than in proportion.} From 16 to 1,024 animals
there are 64 times as many facts. Taking the middle one of the three figures in each entry of
table \ref{tab:effort1}, width 8 needed 8 times as many sets (512 against 64), and widths 16
and 32 needed about 6 times as many.

\textbf{The smaller the model, the more sets it needed, and the less dependable it was.} On
1,024 animals, width 8 needed four times as many sets as width 32 (512 against 128, the middle
figures), and its
three repeats differed widely: 256, 512 and 1,639 sets. The slowest only just made it: it
passed for the second time at the grading made when its 30 minutes ran out. For its first 250 sets its score stayed near zero. Later it stood still
again, at a score of 25 from set 512 to set 724, before climbing to 52 and then to 100. (A
score of 25 is what a model gets if it knows one of the four kinds of fact for every animal, so
it seems to have picked up one kind of fact at a time.) This is a warning for everything that
follows: a model can stand still for a long time and then move again.

